\chapter{Ag Guvenligi}

\section*{Giris}
Ag guvenligi, kurumsal iletisim omurgasini korur. Perimeter odakli yaklasimdan, gorunurluk ve surekli dogrulama odakli mimarilere gecis zorunlu hale gelmistir.

\section{Ag Iletisim Temelleri ve Iletisim Protokolleri}
\subsection{OSI Referans Modeli ve TCP/IP Protokol Yigini}
OSI modeli katmansal dusunceyi, TCP/IP ise pratik internet iletisimini temsil eder. Guvenlik kontrolleri her katmanda farkli uygulanir.

\subsection{TCP 3'lu El Sikisma (3-Way Handshake) ve Oturum Yonetimi}
SYN, SYN-ACK, ACK akisi; oturum acilisi ve anomalilerin tespiti acisindan SOC analizlerinde temel gozlem noktasidir.

\subsection{Yonlendirme Protokolleri Dinamikleri (BGP, OSPF, RIP ve Link State/Distance Vector)}
Yonlendirme guvenligi icin route filtreleme, kimlik dogrulama ve anomali izleme (route hijack tespiti) uygulanmalidir.

\section{Kurumsal Ag Guvenligi Mimarisi}
\subsection{KOBI'lerden Dagitik Isletmelere Ag Tasarim Prensipleri}
Kurum olcegine gore segmentasyon, merkezi loglama ve uzaktan erisim guvenligi kademeli olarak olgunlastirilmalidir.

\subsection{Ag Segmentasyonu ve VLAN Izolasyonu}
VLAN segmentasyonu ve mikro-segmentasyon, yanal hareketin sinirlandirilmasinda temel savunmadir.

\section{Yeni Nesil Guvenlik Cozumleri}
\subsection{Yeni Nesil Guvenlik Duvarlari (NGFW) ve Uretici Pratikleri (Palo Alto, Fortinet)}
Uygulama-farkindalik, SSL inspection ve tehdit istihbarati beslemeleri ile birlikte calistirildiginda NGFW etkisi artar.

\subsection{Saldiri Tespit ve Onleme Sistemleri (IDS/IPS)}
Imza tabanli ve davranissal tespit birlikte kullanilarak hem bilinen hem bilinmeyen saldiri kaliplari yakalanir.

\section{Gelismis Ag Saldiri Vektorleri ve Savunma}
\subsection{ARP Zehirlenmesi (Poisoning) ve MAC Adresi Sahteciligi (Spoofing)}
\subsection{Paket Parcalanma (Fragmentation Overlap) ve Kaynak Tuketimi (SYN Flood, Smurf) Saldirilari}
\subsection{Oturum Calma (Session Hijacking) ve Ortadaki Adam (MitM) Saldirilari}
Bu saldirilara karsi port security, DHCP snooping, DAI, rate-limit ve TLS/mTLS zorunlulugu gibi kontroller uygulanmalidir.

\section{Ag Gorunurlugu (Network Visibility) ve Analiz}
\subsection{Derin Paket Incelemesi (DPI) ve Sifreli Trafik Analizi}
\subsection{Ag Sizma Testleri ve Zafiyet Tarama Metodolojileri}
Wireshark, Zeek ve netflow/sflow telemetrisi; olay tespiti ve adli analizde birlikte degerlendirilmelidir.

\section{Kablosuz Ag ve Uzaktan Erisim Guvenligi}
\subsection{Mesh Aglar, Tailscale ve ZTNA Entegrasyonlari}
Uzaktan erisimde VPN tek basina yeterli degildir; kimlik dogrulama, cihaz posture ve uygulama tabanli erisim modeli birlestirilmelidir.
